% MT67000C
% Hitzschlag, Sonnenstich
% 14.0
% 2013-11-30
\label{m:sonnenstich}
\FormatTable

\begin{tabularx}{\linewidth}{XX}
\noindent\TabHeading{Hitzschlag}
&
\noindent\TabHeading{Sonnenstich}
\\\addlinespace\midrule\addlinespace
% \begin{minipage}{\linewidth}
\begin{itemize}
  \item \TxMassVitMK
  \item Beengende Kleidungsstücke lockern
  \item Flachlagerung an kühlem Ort (Schatten), Beine und Kopf hochlagern. 
  \item Kühlung von \E{außen} (Fächer, mit Eiswürfeln abreiben etc.)
  \item Kühlung von \E{innen}: für \EE{ansprechbare} Patienten viel kühle
  (alkoholfreie) Flüssigkeit, evtl. Elektrolyt-Getränke
  \item Neurocheck inkl. BZ-Messung
\end{itemize}
% \end{minipage}
&
% \begin{minipage}{\linewidth}
\begin{itemize}
  \item Vitale Bedrohung einschätzen
  \TxBeiVit \TxMassVitMK
  \item Flachlagerung an kühlem Ort (Schatten), Kopf hochlagern
  \item Beengende Kleidungsstücke lockern
  \item Kühlung von außen (Fächer, mit Eiswürfeln abreiben, kalte
  Umschläge auf die Stirn etc.)
  \item Kühlung von innen: für \EE{ansprechbare} Patienten kühle
  (alkoholfreie) Flüssigkeit, evtl. Elektrolyt-Getränke
%   \item Schocklagerung u. -bekämpfung. % REVIEW 2010-05-19
  \item Neurocheck inkl. BZ-Messung
\end{itemize}
% \end{minipage}
\\\addlinespace\bottomrule
\end{tabularx}